
\chapter{Esercizi}

\section{Polinomi}
Calcola il resto $R(x)$ della seguente divisione polinomiale in $\mathbb{Z}^2$. 

$$
\frac{(x^5 + x + 1) (x^4 + x^2 + 1) + (x^2+1)}{(x^2 + x)}
$$

Facilitiamo l'esercizio suddividendolo in singoli polinomi.
\begin{itemize}
	\item $A(x) = (x^5 + x^1 + 1)$;
	\item $B(x) = (x^4 + x^2 + 1)$;
	\item $C(x) = (x^2 + 1)$;
	\item $D(x) = (x^2 + x)$;
\end{itemize}

Svolgiamo le operazioni.
\begin{enumerate}
	\item $A(x) \cdot B(x) = (x^9 + \cancel{x^5} + x^4 + x^7 + x^3 + x^2 + \cancel{x^5} + x + 1)$
	
	\item $A(x) \cdot B(x) + C(x) = (x^9 + x^4 + x^7 + x^3 + \cancel{x^2} + x + \cancel{1}) +  (\cancel{x^2} +\cancel{1}) = (x^9 + x^7+ x^4 + x^3 + x)$. 
	
	Abbiamo letteralmente effettuato uno $XOR$. Chiameremo $E(x)$ il polinomio risultante.
	
	\item $E(x)/D(x)$. Facciamo delle osservazioni.
	\begin{itemize}
		\item $E(x) = \verb|1010011010|$
		\item $D(x) = \verb|110|$
	\end{itemize}
	
	Facciamo la divisione.
	
	\begin{verbatim}
		1010011010  |
		110         |
		-------------
		11         |
		110        |
		-------------
		110    |
		110    |
		------------
		10  | resto trovato
	\end{verbatim}
	
	$R(x) = \verb|10| = 1\cdot x^1 + 0\cdot x^0 = x$;
\end{enumerate}

Soluzione: $R(x) = x$.